%%%%

%%%   Самарский государственный технический университет

%%%   Кафедра прикладной математики и информатики

%%%

%%%   Преамбула для статьи в журнал "Вестник Самарского государственного

%%%   технического университета. Серия: «Физико-математические науки»"

%%%   Ориентирована под MikTEx 2.4 for Win32

%%%

%%%   Хотя сам журнал собирается под GNU/Linux на дистрибутиве TeXLive



\documentclass[11pt,a4paper,twoside]{article}



\usepackage[cp1251]{inputenc}

\usepackage[T2A]{fontenc}

\usepackage[english,russian]{babel}

\usepackage{samgtubib}

\usepackage[dvips]{graphicx}

\usepackage[multiple]{footmisc}



\usepackage{samgtu}

\renewcommand{\VestnikYear}{2009} % Год издания Вестника

\renewcommand{\VestnikNumber}{2\,(19)} % Номер Вестника

\renewcommand{\VestnikMonth}{Ноябрь} % Месяц выхода Вестника

\renewcommand{\VestnikData}{01.11.09} % Дата подписания Вестника в печать

\renewcommand{\VestnikUPL}{26,64} % Усл. печ. л.

\renewcommand{\VestnikUIL}{25,73} % Уч.-изд. л.

\renewcommand{\VestnikRegNumber}{479/09} % Регистрационный номер

\renewcommand{\VestnikOrder}{994} % Номер заказа







%%\samgtupubl % для генерации pdf, отдаваемого в типографию СамГТУ

%\pdfcrop % обрезка pdf для размещения на math-net.ru

%\draft % На корректуре

\filllastpage % Последняя страница не заполнена не полностью

%\briefcomm % Статья является кратким сообщением







\vsgtu{1222} % Номер статьи по базе Math-Net.ru



\FirstPage{125}

\LastPage{129}





%\begin{document}



%

\renewcommand{\baselinestretch}{0.95}



%\Partition{Краткие сообщения}{Brief Communications}

%

%\ShortPartition{Математическое моделирование}{}





\udk{517.958:621.785.54/.56}

\msc{76N15,  76L05}



\rutitle{Моделирование процесса увлечения частиц порошка взрывными ударными волнами} {Моделирование процесса увлечения частиц порошка взрывными ударными волнами}



\ruauthor{А.~И.~Крестелев}{А.~И.~К\,р\,е\,с\,т\,е\,л\,е\,в}



\ruaffil{\rusamgtu.}



\email{1}{\mem{a.krestelev@yandex.ru}}



\ruauthordetails{1}{\fullauthor{\rupid{38596}{Анатолий Иванович Крестелев}}\regalia{к.ф.-м.н., доц.}\position{доцент}\dept{каф. общей физики и физики нефтегазового производства}}



\rukeywords{взрывная ударная волна, продукты детонации, частицы порошка}



\ruabstract{На основе простой физической модели рассматривается процесс взаимодействия продуктов детонации во взрывной волне с частицами порошка при взрывном напылении износостойких покрытий.

Увлечение частиц происходит в результате неупругого соударения молекул продуктов детонации с частицами порошка. 

Получено уравнение для определения скорости частиц во фронте волны, образовавшейся при взрыве сферического заряда взрывчатого вещества. 

Анализируются алгоритмы решения полученного уравнения в зависимости от динамических характеристик 

продуктов детонации. 

Полученные результаты могут быть использованы при проектировании технологических схем взрывного напыления покрытий и теоретическом анализе процесса сверхглубокого проникания частиц порошка в~металлические мишени.

}



\entitle{Simulation of the Process of Entrainment of~a~Powder Particles by Explosive Shock Waves}



\enauthor{A.~I.~Krestelev}{A.~I.~K\,r\,e\,s\,t\,e\,l\,e\,v}





\enaffil{\ensamgtu.}



\enauthordetails{1}{\fullauthor{\enpid{38596}{Anatoliy I. Krestelev}} \regalia{Cand. Phys. \& Math. Sci.} \position{Assotiate Professor} \dept{Dept. of General Physics and Physics of Oil and Gas Production}}



\enabstract{The coating on the surface of metals and alloys is used to increase strength and durability of materials. There are many different technological schemes of this process. It is of interest to use the explosion energy to create a stream of particles deposited on the surface of metals. The article is based on a simple physical model describing the process of interaction between the products of detonation of the explosive substance with the particles of a powder in an explosive spraying of wear-resistant coatings. The entrainment of particles occurs due to inelastic collisions of molecules of detonation products with particles of a powder. An equation for determining the velocity of particles in the wave front, formed during the explosion of a spherical explosive charge, is received. The equation of motion of the particles can be written for the case of an explosion of a cylindrical charge of explosive. The solving algorithms of the obtained equation are analyzed depending on the dynamic characteristics of detonation products. The obtained results can be used for designing the technological schemes of explosive deposition and making the theoretical analysis of the process of superdeep penetration of particles of a powder in a metallic target.}



\enkeywords{explosion shock wave, detonation products, powder particles}



\date{19/IV/2013}

\revisiondate{05/XII/2013}

\accepteddate{17/I/2014}



\makerutitle





Исследование взаимодействия потока высокоэнергетических частиц с~поверхностью материалов представляет интерес при решении как  прикладных технологических задач, так и фундаментальных проблем. 

К~первому кругу задач можно отнести разработку технологий нанесения износостойких покрытий на

поверхность конструкционных материалов с помощью взрывного напыления~\cite{krest:bib1,krest:bib2},

с~другой стороны, очень важным является изучение самого процесса взаимодействия взрывных ударных волн с частицами распыляемого порошка.



В настоящей работе с помощью достаточно простой физической модели исследуется движение частиц порошка, инициируемое продуктами детонации (ПД) взрывной ударной волны в~воздухе. 

В~зависимости от технологической схемы взрывного напыления порошковая навеска либо непосредственно примыкает к~корпусу заряда взрывчатого вещества (ВВ), либо находится вблизи заряда при канальной схеме устройства~\cite{krest:bib1}. 

В~любом случае частицы порошка увлекаются расширяющимися в окружающее пространство продуктами детонации взрывной ударной волны. 

Механизм взаимодействия частиц порошка с молекулами ПД можно смоделировать следующим образом.



Частица порошка начинает двигаться в результате неупругого соударения с~молекулами продуктов детонации, расширяющимися в окружающее пространство.

Процесс последовательных столкновений частицы с молекулами ПД можно описать с~помощью системы уравнений, представляющей собой цепочку законов сохранения импульса:

\begin{equation}

\label{krest:eq1}

\begin{aligned}

& mu+M v_0 =(M+m) v_1, \\

& mu+(M+m) v_1 =(M+2 m) v_2, \\

& \cdots \\

& m u+[M+(N-1) m]  v_{N-1} =(M+N m) v_N,

\end{aligned}

\end{equation}

где $m$ "--- масса молекул продуктов детонации, $u$ "--- массовая скорость ПД, $M$ "---

масса частицы порошка, $v_{0}$ "--- начальная скорость частицы порошка, $v_{N}$ "--- скорость после $N$ соударений.



В рамках исследуемой модели массы всех частиц продуктов детонации считаются

одинаковыми. Если учесть, что в состав ПД входят такие молекулы, как O$_2$,

N$_2$, CO, H$_2$O, то это допущение можно считать оправданным. 

Кроме того, при решении системы учитывалось, что $M\gg m$. 

Если учесть, что масса частицы $M$ значительно превосходит массу молекул продуктов детонации $m$, то из системы \eqref{krest:eq1} следует, что скорость частицы порошка после $N$ соударений 

\begin{equation}

\label{krest:eq2}

v_N =N \frac{m}{M} u+v_0.

\end{equation}



В дальнейших расчётах будем считать начальную скорость частицы равной нулю. 

Для определения числа соударений $N$ будем считать, что за время $d\tau$ с частицей

порошка ударятся все молекулы ПД, находящиеся в объёме~$dV =S (u-v) d\tau$, где $S$ "--- площадь поперечного сечения частицы. 

Здесь учтено, что частица порошка также движется со скоростью $v$. 

Если концентрация продуктов детонации равна $n_d$, то за время $t$ произойдет $N$ соударений:

\begin{equation}

\label{krest:eq3}

N=\int_0^t n_d S  (u-v)\, d\tau.

\end{equation}



Из соотношения \eqref{krest:eq3} следует, что число столкновений частицы является функцией

скорости её движения, а следовательно, уравнение \eqref{krest:eq2} преобразуется в~интегральное уравнение вида

\begin{equation}

\label{krest:eq4}

v=\frac{m}{M} u \int_0^t n_d S  (u-v)\,d\tau.

\end{equation}



В теории взрывных ударных волн чаще используется не концентрация продуктов детонации, а их плотность $\rho$. 

Если учесть, что $n_d =\rho/m$, то уравнение \eqref{krest:eq4} запишется в виде

\begin{equation}

\label{krest:eq5}

v+\frac{u}{M} \int _0^t \rho S v\, d\tau =\frac{u}{M} \int _0^t \rho S u\, d\tau.

\end{equation}



В результате расширения продуктов детонации в окружающую среду меняются и их

параметры $\rho $ и $u$. 

Плотность ПД можно взять среднюю по объёму, тогда $\rho = m_0/V$, где $m_0$ "--- масса заряда взрывчатого вещества. 

Тогда для сферического заряда ВВ плотность составит величину

\begin{equation}

\label{krest:eq6}

\rho =\frac{3 m_0}{4 \pi r^3},

\end{equation}

где $r$ "--- расстояние до фронта ударной волны в воздухе.



Для цилиндрического заряда ВВ имеем

\begin{equation}

\label{krest:eq7}

\rho =\frac{m_0}{\pi l r^2},

\end{equation}

где $l$ "--- длина цилиндрического заряда, при этом предполагается, что расширение ПД происходит только в~радиальном направлении. 

Это справедливо в~том случае, когда длина заряда значительно больше его диаметра.



Если учесть, что плотность и другие параметры продуктов детонации являются функциями расстояния~\cite{krest:bib3}, то в~уравнении \eqref{krest:eq5} удобно перейти от переменной

времени к~переменной расстояния от центра заряда ВВ до фронта ударной волны расширяющихся продуктов детонации: 

$$dr=v_d  \,  d\tau ,$$ где 

$v_d $ "--- скорость фронта ударной волны. 

Скорость ударной волны в воздухе является переменной величиной~\cite{krest:bib4}. 

При переходе к пространственной переменной уравнение \eqref{krest:eq5} для сферического заряда ВВ запишется в виде

\begin{equation}

\label{krest:eq8}

v+A u  \int_{r_0}^r \frac{v}{r^3 v_d }dr = A u \int_{r_0}^r \frac{u}{r^3  v_d}dr,

\end{equation}

где $A=(3 m_0 S)/(4 \pi M)$, $r_0$ "--- расстояние от центра заряда взрывчатого вещества до частицы порошка в~начальный момент времени. 

Если порошок закреплен на поверхности заряда ВВ, то $r_0$ "--- радиус самого заряда.



Уравнение \eqref{krest:eq8} является интегральным уравнением типа Вольтерра и позволяет определить скорость движения частицы порошка под действием ПД взрывной ударной волны, если известны законы изменения массовой скорости ПД и скорости фронта ударной волны.



Массовую скорость продуктов детонации $u$ при истечении их в воздух можно определить через давление ПД соотношением~\cite{krest:bib4}

\begin{equation}

\label{krest:eq9}

u=\frac{D}{n+1}\Bigl[

1+\frac{2 n}{n-1} (1-N_{KH} )+\frac{2 n N_{KH}}{k-1} 

\left( 1-\left( {P}/{P_K } \right)^{(k-1)/(2k)}\right) \Bigr],

\end{equation}

где $N_{KH} =\left({P_K }/{P_H }\right)^{({n-1})/(2n)}$,

$P_K$, $P_H$ "--- давления продуктов детонации в точке сопряжения и точке Чепмена"--~Жуге соответственно~\cite{krest:bib3}, $n=3$ "--- показатель изоэнтропы продуктов детонации в зоне химической реакции.



Массовая скорость ПД \eqref{krest:eq9} входит в интегральное уравнение \eqref{krest:eq8} и существенно усложняет его,

поэтому можно использовать более простую связь массовой скорости и давления продуктов детонации~\cite{krest:bib4}:

\begin{equation}

\label{krest:eq10}

u=\sqrt {\frac{2}{k+1} \frac{P}{\rho}},

\end{equation}

где $k\approx 1.2$ "--- показатель изоэнтропы воздуха, 

$P$ "--- давление в волне разрежения продуктов детонации, 

$\rho$ "--- плотность ПД. 

Изменение давления продуктов детонации при расширении их в газовой среде анализировалось в~работе~\cite{krest:bib3}, в которой были получены соотношения, определяющие зависимость давления от

расстояния до фронта ударной волны. 

Например, в случае сферического заряда ВВ давление определяется выражением

\begin{equation}

\label{krest:eq11}

P =

\left\{

\begin{array}{ll}

 \alpha P_H ( {r}/{r_0} )^{-9}, & P_K < P\le P_H, \\

 \beta P_K ( {r}/{r_0} )^{-3 k}, & P\le P_K,

\end{array}

\right.

\end{equation}

где $r_0$ "--- радиус заряда ВВ. 

Соотношение \eqref{krest:eq11} получено в предположении о~ступенчатом изменении показателя изоэнтропы в~точке сопряжения.



Таким образом, в работе получено уравнение для определения скорости движения частиц порошка при взрывном напылении его на поверхность материалов. 

Динамические характеристики частиц, увлекаемых взрывной ударной волной, могут представлять интерес и при исследовании такого явления, как сверхглубокое проникание их в металлы, которое экспериментально наблюдалось и теоретически исследовалось многими авторами~\cite{krest:bib5,krest:bib6}.



\begin{thebibliography}{99}



\RBibitem{krest:bib1}

\jour Деформация и разрушение

\yr 2008

\issue 5

\pages 44--47

\by М.~К.~Валюженич, А.~Л.~Кривченко, А.~М.~Штеренберг

\paper Модификация поверхности титановых сплавов взрывоплазменным напылением

\elib{http://elibrary.ru/item.asp?id=11969177}

\misctransl

\by M.~K.~Valyuzhenich, A.~L.~Krivchenko, A.~M.~Shterenberg

\paper Modification of the surface of titanium alloy by explosion plasma spraying  coatings

\jour Deformatsiya i razrusheniye materialov

\yr 2008

\issue 5

\pages 44–47

\lang In Russian





\RBibitem{krest:bib2}

\by А.~И.~Крестелев, М.~К.~Валюженич, Л.~А.~Довбня

\paper Использование взрывных ударных волн для нанесения покрытий на поверхности металлов

\inbook Сб. тезисов XVII Международной конференции ``Физика прочности и пластичности материалов''

\publaddr Самара

\publ СамГТУ

\yr 2009

\pages 239--240

\misctransl

\by A.~I.~Krestelev, M.~K.~Valyuzhenich, L.~A.~Dovbnya

\paper The use of explosive shock waves for coating on metal surfaces

\inbook Sb. tezisov XVII Mezhdunarodnoy konferentsii ``Fizika prochnosti i plastichnosti materialov''

\rm [Abst. XVIIth Int. Conf. ``Physics of Strength and Plasticity of Materials'']

\publ Samara State Technical Univ.

\publaddr Samara

\yr 2009

\pages 239--240

\lang In Russian



\RBibitem{krest:bib3}

\by А.~И.~Крестелев

\paper Термодинамический анализ взаимодействия взрывных ударных волн с поверхностью материалов

\jour Известия Самарского научного центра Российской академии наук

\yr 2011

\vol 13

\issue 6

\pages 72--76

\elib{http://elibrary.ru/item.asp?id=17846092}

\misctransl

\by A.~I.~Krestelev

\paper Thermodynamic analysis of interaction of explosive shock waves with the surface of materials

\jour Izvestiya Samarskogo nauchnogo tsentra Rossiyskoy akademii nauk

\yr 2011

\vol 13

\issue 6

\pages 72--76

\lang In Russian







\RBibitem{krest:bib4}

\book Физика взрыва

\bookvol 1

\ed Л.~П.~Орленко

\publaddr М.

\publ Физматлит

\yr 2004

\totalpages 823

\misctransl

\ed L.~P.~Orlenko

\book Fizika vzryva \rm [Physics of explosion]

\bookvol 1

\publaddr Moscow

\publ Fizmatlit

\yr 2004

\totalpages 823

\lang In Russian





\RBibitem{krest:bib5}

\by Л.~В.~Альтшулер, С.~К.~Андиленко, К.~Г.~Романов, С.~М.~Ушеренко 

\paper О модели сверхглубокого проникания

\jour Письма в журнал технической физики

\yr 1989

\vol 15

\issue 5

\pages 55--57

\misctransl

\by L.~V.~Al'tshuler, S.~K.~Andilenko, K.~G.~Romanov, S.~M.~Usherenko

\paper On the model of the super deep penetration

\jour Pis'ma v zhurnal tekhnicheskoy fiziki

\yr 1989

\vol 15

\issue 5

\pages 55--57

\lang In Russian



\RBibitem{krest:bib6}

\by О.~В.~Роман, С.~К.~Андилевко, С.~С.~Карпенко

\paper Некоторые параметры сверхглубокого проникания порошка в алюминиевую преграду

\jour Химическая физика

\yr 2002

\vol 21

\issue 9

\pages 52--56

\misctransl

\by  O.~V.~Roman, S.~K.~Andilevko, S.~S.~Karpenko

\paper Some parameters of super deep penetration of powders of bronze into aluminum obstacle

\jour Khimicheskaya Fizika

\yr 2002

\vol 21

\issue 9

\pages 52--56

\lang In Russian





\end{thebibliography}



%\shortengtitle{A.~I.~K\,r\,e\,s\,t\,e\,l\,e\,v}{Simulation of the Process of Entrainment of Particles \dots}



\renewcommand{\baselinestretch}{0.9}



\makeentitle



\newpage





%\end{document}

